\section{Concluding remarks}\label{sec:conclusio}
\rev{In a globalised world, monetary policy depends not only on prevailing macroeconomic fundamentals but also on how central banks interact, because key transmission channels have an explicit international dimension and effective policy hinges on international spillovers in financial and real markets as much as on domestic first-round effects. In this paper, we develop an empirical weekly time series model that captures, in a nonlinear fashion, the interactions between domestic macro-financial dynamics in the euro area (EA) and the US monetary policy stance. We identify shocks using readily available high-frequency instruments designed to capture unexpected movements in monetary policy, and we measure the US policy stance from an estimate of the neutral interest rate.}

\rev{In sum, our results reveal a systematic interaction between US and EA monetary policy. Rather than relying on a comparison of two extreme stances, we characterise the transmission of an ECB surprise across the entire range of the Fed's stance. For nominal bond yields this mapping is clean and monotone: a contractionary EA surprise raises euro area yields significantly when the Fed is accommodative, and this reaction weakens steadily as the Fed turns restrictive, vanishing or reversing at long maturities under a strongly hawkish Fed.} These effects reverse after about two weeks. At this horizon, bond yields tend to fall in response to capital inflows into the EA economy, which pushes up bond prices. In contrast, when the EA shock occurs under a restrictive Fed stance, we observe somewhat weaker reactions with heterogeneity across maturities. Looking at the reactions of the inflation swap curve, we observe muted differences when controlling for the US stance. In general, the results point to negative reactions at impact, with different dynamics thereafter depending on the prevailing US monetary policy stance. However, even though they are marginally significant, there is some evidence of nonlinearities.
\rev{Finally, our model lets us examine real interest rate reactions by taking the difference between the nominal-yield and inflation-swap IRFs; these suggest that real rates rise in all scenarios, though more so when EA shocks occur during a restrictive Fed stance.}

\rev{The revision of this paper has substantially broadened the evidence behind these conclusions. The state dependence with respect to the Fed's stance survives the inclusion of global variables (oil, the broad dollar, and global risk aversion) and the use of a stance measure purged of the global cycle; it is empirically distinct from state dependence on the euro area business cycle, as the regime allocations implied by the two conditioning variables are essentially uncorrelated; it is present for both tightening and easing surprises; and it is estimated on a sample running through the end of 2025 that covers the post-pandemic global tightening, during which the model identifies a pronounced restrictive Fed regime for 2022--2024 that sharpens the identification of the state dependence. Moreover, since the model is specified in first differences, the fast mean reversion of the reported responses translates into persistent level effects on yields, with the sign of the long-run level effect itself depending on the Fed's stance.}

\rev{The model interacts foreign monetary policy with domestic quantities in a simple and transparent way. A natural next step is to relax the logistic form of the interaction and let the data determine how the transmission depends on the Fed's stance --- for instance through a time-varying parameter VAR whose coefficients are nonparametric functions of the stance, modelled with regression trees \citep{hauzenberger2025bayesian}. Such an approach would accommodate richer, potentially non-monotone patterns of state dependence than the two-regime mixture imposed here.} In addition, we have only included EA-based quantities and capture international spillovers mainly by including the US monetary policy stance. Estimating larger bi-country models would in principle be feasible, which we leave for further research.

